\section{The integral matrix scheme and the later formal constructions}

\begin{statement}{B.1}{Proved}{Equations and a scheme-theoretic factorization}
Let

$$
H=\Spec
 \frac{\Z[a,b,c,d,e,t,(c^2+d^2)^{-1}]}
 {(e^2-1,\ t^2-1,\ a^2+ab+b^2-et(c^2+d^2))}.
$$

Then $M$ is affine and

$$
M\simeq H\times\SL_2
$$

as schemes. The isomorphism sends $(h,g)$ to
$(x(h),y(h),g\operatorname{diag}(eq,1))$.
In particular the determinant relation in C3-399 is
$\det f=\det y=t\det x$, not equality of three unrelated
norm-one conditions.
\end{statement}
\begin{proof}
The equations defining $H$ are exactly the first two matrix
coordinates and their determinant compatibility. Since
$e,t,q$ are units, so is $n=etq$.
A matrix $f$ of determinant $eq$ factors uniquely as
$g\operatorname{diag}(eq,1)$ with $g\in\SL_2$.
The factorization and its inverse are regular over the displayed
coordinate ring. This proves the scheme assertion over $\Z$,
including all nonreduced test rings.
\end{proof}

\begin{statement}{B.2}{Proved}{Integral flatness and dimension}
The group scheme $M$ is flat and finitely presented over $\Z$,
of relative dimension six. The three defining polynomials of
$H$, before localization, form a monic Groebner basis with
leading monomials $e^2,t^2,a^2$ for a suitable lexicographic
order. This is an integral certificate, not a computation
over $\Q$ followed by clearing denominators.
\end{statement}
\begin{proof}
Take the variable order $a>e>t>b>c>d$.
The three leading coefficients are one and the leading
monomials are pairwise relatively prime. The product criterion,
or direct division of their three S-polynomials, gives zero
remainders over $\Z$. The residue monomials have exponents
at most one in $a,e,t$, with arbitrary exponents in $b,c,d$.
They form a free $\Z$-basis. Localizing at $q=c^2+d^2$
preserves $\Z$-flatness.
Every fibre has dimension three: the two sign coordinates
are finite, and the monic equation in $a$ leaves a finite
algebra over the three polynomial coordinates $b,c,d$.
The nonempty unit locus still has dimension three.
Multiplication by the smooth three-dimensional scheme $\SL_2$
in B.1 proves the assertion for $M$.
\end{proof}

\begin{statement}{B.3}{Proved}{The smoothness locus and the connected kernel}
The scheme $M$ is smooth exactly over $\Spec\Z[1/2]$.
The subgroup $M^{00}=\ker(e,t)$ is smooth over all of $\Z$,
with geometrically connected fibres. It is therefore a
normal integral scheme. The map

$$
M\longrightarrow V=\mu_2^2,\qquad m\longmapsto(e,t)
$$

is an fppf surjection with kernel $M^{00}$.
\end{statement}
\begin{proof}
Away from two the signs are etale. In the remaining equation,
the derivatives in $c,d$ are $-2etc,-2etd$.
They cannot vanish simultaneously on $q$ invertible.
The standard smoothness criterion therefore applies
\cite[Lemma10.137.6]{AB_Stacks}.

For $M^{00}$ the equation is $n=q$.
In characteristic two, its derivatives in $a,b$ are $b,a$,
and these cannot both vanish because $n=q$ is a unit.
In characteristic three, simultaneous vanishing of
$2a+b,a+2b$ forces $a=b$, hence $n=0$, again impossible.
In every other characteristic, the same conclusion follows
from invertibility of three, or from the derivatives in $c,d$.
Thus $M^{00}$ is smooth over $\Z$.

Over an algebraically closed field of characteristic different
from two and three, its scalar factor is
$x_1x_2=y_1y_2$ with all four variables units, hence a
three-dimensional torus. In characteristic three it is
$\Ga\times\Gm^2$: the Eisenstein norm is the square of a
linear coordinate, while the Gaussian algebra is split.
In characteristic two it is also $\Ga\times\Gm^2$ after
splitting the etale Eisenstein algebra.
Each description is connected, as is $\SL_2$.
Smoothness over the regular base and the integral generic
fibre imply normality and integrality of the total scheme.

For any prescribed signs, fppf locally choose square roots
of their required unit norms; scalar matrices then supply
$x,y$, and choose a diagonal $f$ of the resulting determinant.
This proves surjectivity onto $V$. Each fibre is a torsor
under the smooth kernel, so the quotient map is smooth
and fppf. B.4 shows that $M$ itself is not smooth at two.
\end{proof}

\begin{statement}{B.4}{Proved}{The characteristic-two fibre}
The geometric characteristic-two fibre of $M$ has one
irreducible component, dimension six, multiplicity four,
and tangent dimension eight at every geometric point.
Its reduced group is

$$
\SL_2\times ST_r\times\Gm\times\Ga
$$

over $\F_2$. The whole characteristic-two fibre is singular
as a six-dimensional fibre scheme.
\end{statement}
\begin{proof}
In characteristic two,
$e^2-1=(e-1)^2$ and $t^2-1=(t-1)^2$.
Their differentials vanish. The norm equation has nonzero
differential in the $a,b$ directions on the unit locus.
Hence the scalar tangent dimension is five, rather than
three, and the $\SL_2$ factor adds three.
The sign base $\mu_2^2$ has local length four.
The smooth quotient map of B.3 shows that the fibre is a
smooth deformation over this length-four local base of its
reduction; equivalently its completed local rings have two
independent square-zero sign parameters and six smooth
parameters. This proves the multiplicity.

On the reduction $e=t=1$. Put $z=c+d$, a unit, and write
$x=z x_0$. Then $N_r(x_0)=1$, while $d/z$ is an arbitrary
additive coordinate. These substitutions and their inverses
give the displayed reduced group scheme.
The Eisenstein norm-one group is a connected torus at two,
so the reduced fibre is geometrically irreducible.
The tangent computation proves singularity at every point.
\end{proof}

\begin{statement}{B.5}{Proved}{The characteristic-three fibre}
The geometric characteristic-three fibre of $M$ has four
smooth irreducible components. Each is isomorphic as a
scheme to

$$
\SL_2\times\Ga\times\Gm^2.
$$

Thus three is not a bad smoothness prime for this particular
six-dimensional scheme.
\end{statement}
\begin{proof}
The two signs are etale and give four components.
In characteristic three, $n=(a-b)^2$.
After extending the field to split the Gaussian algebra,
write $q=uv$, with $u,v$ units.
For fixed signs the equation is $z^2=etuv$, where
$z=a-b$ is a unit. It eliminates $v$ as $et z^2/u$;
$b$ is free. This gives $\Ga\times\Gm^2$.
The factorization in B.1 completes the proof.
\end{proof}

\begin{statement}{B.6}{Proved}{The norm-one groups have different bad fibres}
The group $ST_r$ is smooth over all of $\Z$.
Its geometric fibre at three is the disjoint union of two
copies of $\Ga$. The group $ST_s$ is nonreduced at two:
there its equation is $(c+d-1)^2=0$.
Neither fact contradicts smoothness of $M^{00}$.
\end{statement}
\begin{proof}
For $ST_r$, the equation is $a^2+ab+b^2=1$.
The gradient analysis in B.3 proves smoothness at every prime.
At three it becomes $(a-b)^2=1$, giving the two affine lines.
For $ST_s$ the equation is $c^2+d^2=1$, which at two has
the displayed square form. Away from two its gradient is
nonzero. In $M^{00}$, the Gaussian norm is equated to an
independent Eisenstein norm rather than to a fixed constant;
the latter provides the missing derivative at two.
\end{proof}

\begin{statement}{B.7}{Proved}{The main quotient and its exact kernel}
The homomorphism

$$
\Psi:M\longrightarrow\GL_2\times\mu_2^2,\qquad
 (x,y,f)\longmapsto(f,e,t)
$$

is fppf surjective, with kernel
$K=ST_r\times ST_s$.
This kernel is central in $M^{00}$; it is not central in
the full group $M$.
\end{statement}
\begin{proof}
For target data $(f,e,t)$, put $\delta=\det f$.
One needs $N_r(x)=t\delta$ and $N_s(c-di)=e\delta$.
Adjoining square roots of these two units gives finite
locally free covers on which scalar solutions exist.
This proves fppf surjectivity. Setting $f=I,e=t=1$
gives exactly the two norm-one factors.
Inside $M^{00}$ both centralizer factors are commutative.
The $e=-1$ component conjugates the Gaussian factor by its
nontrivial algebra involution, which is inversion on $ST_s$.
That action is nontrivial, so centrality fails in $M$.
\end{proof}

\begin{statement}{B.8}{Proved}{The projective kernel and central-product presentation}
Put

$$
H'=\{(x_0,y_0):x_0\in T_r,\ N_r(x_0)^2=1,\
                 y_0\in N'_s,\ \det y_0=1\}.
$$

Then

$$
R_M=\ker(M\to\PGL_2)\simeq\Gm\times H'.
$$

There is an fppf exact central-product sequence

$$
1\longrightarrow\mu_2\longrightarrow R_M\times\SL_2
 \longrightarrow M\longrightarrow1,
$$

where the kernel is the antidiagonal scalar copy of $\mu_2$.
\end{statement}
\begin{proof}
In the projective kernel, $f=\lambda I$ uniquely.
Write $x=\lambda x_0$, $y=\lambda y_0$.
The determinant conditions become
$N_r(x_0)=t\in\mu_2$ and $\det y_0=1$.
This gives the asserted product, with its explicit inverse.
The embedded $\SL_2$ occupies only the third matrix factor
and commutes with $R_M$.
Given $m=(x,y,f)$, fppf locally choose $\lambda^2=\det f$.
Then $(x,y,\lambda I)\in R_M$ and
$\lambda^{-1}f\in\SL_2$, proving surjectivity.
The product is the identity exactly when $x=y=I$ and the
two third factors are inverse scalar matrices with square one.
\end{proof}

\begin{statement}{B.9}{Proved}{The two centres}
The centres as group schemes are

$$
Z(M^{00})=\Gm\times ST_r\times ST_s,
 \qquad
 Z(M)=\Gm\times S'T_r\times\mu_2.
$$

In particular $Z(M^{00})$ is not the scalar-coordinate group
$H^{00}$ from B.1. Over $\Z$ their point counts are
$48$ and $24$, respectively.
\end{statement}
\begin{proof}
Commutation with the embedded $\SL_2$ forces the third
matrix to be scalar, say $\lambda I$.
In $M^{00}$ both other factors commute, and the norm equations
give $x=\lambda x_0$, $y=\lambda y_0$ with both $x_0,y_0$
of norm one. This gives the first formula.

For the full group, a central $y$ must be scalar.
Indeed commutation with the universal Gaussian unit forces
its sign to be one, and commutation with the universal sign
over $R[e]/(e^2-1)$ forces its off-diagonal coefficient to
vanish. This last universal test is necessary: checking only
the two ordinary signs would lose information at two.
Write $y=\lambda\kappa I$. Then $\kappa^2=1$ and
$N_r(x/\lambda)\in\mu_2$, giving the second formula.
Conversely these elements commute with every factor.

The integral norm-one units have orders six and four.
Thus the first centre has $2\cdot6\cdot4=48$ integral points.
In $H^{00}(\Z)$ positivity forces both norms to one,
giving only $6\cdot4=24$ points.
\end{proof}

\begin{statement}{B.10}{Proved}{Graph subgroups and the true line stabilizer}
Let $\chi=\det f=eq$ on $H$.
The graph $f=\operatorname{diag}(\chi,1)$ defines a subgroup
$M(0)\simeq H$. The inverse image $M[0]$ of $B'_0$ has
dimension four. The graph $f=y$ defines $M(s)\simeq H$.
The graph $f=x$ is available only on the locus $t=1$.
The full stabilizer of the first coordinate line under
$M\to\GL_2$ is the inverse image of $B_0$, of dimension five;
it is strictly larger than $M[0]$.
\end{statement}
\begin{proof}
Each displayed graph is a graph of a group homomorphism.
For $f=x$, equality $\det f=t\det x$ forces $t=1$,
because $\det x$ is a unit.
In $B'_0$, the determinant specifies the upper-left entry,
leaving one free upper-right entry over $H$.
The full upper triangular group also has an arbitrary
bottom-right unit, leaving two parameters over $H$.
For example $\operatorname{diag}(2,1/2)$ over $\Q$ fixes
the first coordinate line but does not belong to $B'_0$.
Thus a line-normalizer statement cannot substitute $M[0]$
for the full inverse image of $B_0$.
\end{proof}

\begin{statement}{B.11}{Proved}{The seven-dimensional auxiliary scheme}
The fibre product

$$
SM=M\times_{\chi,\Gm,\det}B'_0
$$

is affine, flat over $\Z$, and has relative dimension seven.
Its projection to $M$ is a split smooth morphism with kernel
$\Ga$. It is not $M^+=\ker\chi$.
\end{statement}
\begin{proof}
The determinant map $B'_0\to\Gm$ is the projection
$(m,n)\mapsto m$ and has the section $m\mapsto(m,0)$.
Its base change has the stated properties and is isomorphic
as a scheme to $M\times\A^1$.
It adds a Borel parameter instead of imposing determinant one.
\end{proof}

\begin{statement}{B.12}{Proved}{Exact finite counts}
The scalar scheme and matrix group have the following point counts:

$$
\begin{array}{c|r|r|r|r}
R&|H(R)|&|\SL_2(R)|&|M(R)|&|M^{00}(R)|\\\hline
\F_2&6&6&36&36\\
\F_3&96&24&2304&576\\
\F_5&384&120&46080&11520\\
\F_7&1152&336&387072&96768\\
\F_{11}&5760&1320&7603200&1900800\\
\F_{13}&6912&2184&15095808&3773952\\
\Z/4&192&48&9216&2304\\
\Z/8&6144&384&2359296&147456\\
\Z/9&2592&648&1679616&419904
\end{array}
$$

These are counts of ring-valued points, not fibre multiplicities.
\end{statement}
\begin{proof}
For each unit $q$, enumerate the Gaussian norm fibre and
the Eisenstein norm fibre at $etq$, and sum over all actual
solutions of $e^2=t^2=1$ in the ring.
For $H^{00}$ use $e=t=1$.
Multiply by $|\SL_2(R)|$ using B.1.
The replacement program enumerates these finite sets exactly.
Its simultaneous triple count from A.8 independently agrees
with the special-linear factor in every row.
\end{proof}

\begin{statement}{B.13}{Proved}{The cusp quotient and its central ambiguity}
Let $C$ be a group, let $L\subseteq C$ be abelian, and let
$B\subseteq\Aut(C)$ preserve $L$.
In $C\rtimes B$ the subgroup

$$
D_L=\{(\ell^{-1},\Int\ell):\ell\in L\}
$$

is normal and centralizes the copy of $C$.
The quotient $G=(C\rtimes B)/D_L$ contains $B$ injectively;
the kernel of $C\to G$ is exactly $L\cap Z(C)$.
Its intersection between the images of $B$ and $C$ is
the image of $L$.

For profinite data with $L,B$ closed and all maps continuous,
these are profinite quotients. If an automorphism of a
profinite group satisfies $u(L)\subseteq L$ for procyclic
$L$, then in fact $u(L)=L$.
\end{statement}
\begin{proof}
Multiplication and conjugation show that the displayed graph
is a subgroup, that it centralizes $C$, and that $B$ conjugates
it into itself. Its intersection with $C$ consists of precisely
the central elements of $L$, whereas its intersection with
$B$ is trivial. The equation
$(\ell,1)=(1,\Int\ell)$ in the quotient proves the intersection
assertion.
For profinite groups the graph is compact and closed.

For the final assertion write $L=\overline{\langle\lambda\rangle}$
and $u(\lambda)=\lambda^a$ for a profinite exponent $a$.
If multiplication by $a$ were not onto $L$, some prime $p$
would divide $a$ on a nontrivial $p$-primary part of $L$.
But
$\lambda=(u^{-n}\lambda)^{a^n}$
for every $n$. In every finite quotient, sufficiently large
$n$ kills the $p$-part of this power. Hence the image of
$\lambda$ has trivial $p$-part in every finite quotient,
a contradiction. Thus multiplication by $a$ is onto $L$.
For finitely generated profinite $C$, characteristic finite
quotients give the usual compact topology on $\Aut(C)$;
the cusp stabilizer is closed. Continuous bijections between
the resulting compact quotients and Hausdorff target groups
are homeomorphisms.
\end{proof}

\begin{statement}{B.14}{Proved}{Point images are norm conditions}
A triple $(f,e,t)$ belongs to $\Psi(M(R))$ if and only if

$$
t\det f\in N_r(T_r(R)),\qquad
 e\det f\in N_s(T_s(R)).
$$

Thus fppf surjectivity in B.7 does not imply pointwise
surjectivity. For the requested small rings the conditions are:

$$
\begin{array}{c|l}
R&\text{condition}\\\hline
\Z&e=t=\det f\in\{1,-1\}\\
\F_2&\text{none}\\
\F_3&t=\det f,\quad e\text{ arbitrary}\\
\Z/4&e=\det f,\quad t\text{ arbitrary}.
\end{array}
$$

\end{statement}
\begin{proof}
Necessity is the defining determinant equation, and independent
solutions of the two norm equations give sufficiency.
Over $\Z$ both unit norms are one by positivity.
Over $\F_3$ the Eisenstein unit norm is one and the Gaussian
norm is onto. Over $\Z/4$ the Gaussian unit norm is one and
the Eisenstein norm is onto the two unit classes.
The finite enumerations in B.12 verify all these norm images.
D.2 gives the complete local and adelic versions.
\end{proof}

\begin{statement}{B.15}{Proved}{The pinning family with a unit parameter}
For a unit $v$, replace the fixed pair by

$$
r_v=\begin{pmatrix}0&1\\v&1\end{pmatrix},\qquad
 s_v=\begin{pmatrix}0&-1\\-v&0\end{pmatrix}.
$$

The two norm polynomials become

$$
n_v=a^2+ab-vb^2,\qquad q_v=c^2-vd^2.
$$

The corresponding matrix group is flat of relative dimension
six, smooth away from two; its kernel with both signs one
is smooth over the entire unit-parameter base.
\end{statement}
\begin{proof}
Direct multiplication gives
$r_v^2=r_v+vI$, $s_v^2=vI$, and $r_v+s_v=E_{22}$.
The centralizer and signed-normalizer matrices have determinants
$n_v$ and $e q_v$. The monic basis argument of B.2 is unchanged.
Away from two the derivatives in $c,d$ cannot vanish on
$q_v$ invertible because $v$ is a unit.
At two the derivatives in $a,b$ are $b,a$ and the unit norm
excludes their simultaneous vanishing.
These identities are included in the universal polynomial
verification program.
\end{proof}

\begin{statement}{B.16}{Proved}{The two adelic projective quotients differ}
The natural map

$$
\GL_2(\widehat\Z)/\{\pm I\}_{\rm diagonal}
 \longrightarrow(\GL_2/\mu_2)(\widehat\Z)
$$

is neither injective nor surjective. Its kernel is

$$
\frac{\prod_p\{\pm1\}}{\{\pm1\}_{\rm diagonal}},
$$

and its cokernel is
$\widehat\Z^\times/\widehat\Z^{\times2}$.
Consequently a requested passage between later profinite groups
cannot treat these two quotients as identical.
\end{statement}
\begin{proof}
Every finite-rank constant-rank projective module over
$\widehat\Z=\prod_p\Z_p$ is free: represent it by an idempotent
matrix, choose a basis and inverse basis on each component,
and assemble those matrices in the product ring.
In particular $\operatorname{Pic}(\widehat\Z)=0$ and all
rank-two vector bundles are trivial.
The Kummer exact sequence and the central quotient sequence give

$$
\mu_2(\widehat\Z)\to\GL_2(\widehat\Z)
 \to(\GL_2/\mu_2)(\widehat\Z)
 \to H^1(\widehat\Z,\mu_2)\to H^1(\widehat\Z,\GL_2).
$$

The last set is trivial and the preceding group is
$\widehat\Z^\times/\widehat\Z^{\times2}$
\cite[Lemma59.28.3]{AB_Stacks}.
Moreover $\mu_2(\widehat\Z)=\prod_p\{\pm1\}$.
Dividing only by the two diagonal signs gives the kernel stated.
Both that kernel and the square-class quotient are nontrivial.
This calculation supplies an obstruction, not a construction
of every marked map between the later tower groups.
\end{proof}

\begin{statement}{B.17}{Proved}{The general reconstruction theorem}
Let $\theta:N\to\Aut(C)$, let $A\triangleleft N$, and let
$\phi:A\to C$ be a homomorphism satisfying

$$
\theta_a=\Int(\phi(a)),\qquad
 \phi(nan^{-1})=\theta_n(\phi(a)).
$$

Then

$$
D_\phi=\{(\phi(a)^{-1},a):a\in A\}
 \triangleleft C\rtimes_\theta N
$$

and $G=(C\rtimes N)/D_\phi$ fits into

$$
1\to C\to G\to N/A\to1.
$$

The map $N\to G$ has kernel $\ker\phi$.
The same assertion holds for profinite groups with closed
continuous data, and for fppf group sheaves as a sheaf quotient.
\end{statement}
\begin{proof}
The inner-action identity shows that the graph is a subgroup
and centralizes $C$. The equivariance identity gives its
normality under $N$. Its intersection with $C$ is trivial,
and its intersection with $N$ is $\ker\phi$.
Every element with second coordinate in $A$ becomes an
element of $C$, proving the exact sequence.
All operations are continuous in the profinite case, and
the graph is closed. The same elementwise identities are
identities of sheaves. They do not on their own prove
representability of a general sheaf quotient.
\end{proof}

\begin{statement}{B.18}{Refuted}{The unrestricted normality assertion}
In a split extension $C\rtimes\Gamma$, the complementary
copy of $\Gamma$ need not be normal.
A counterexample compatible with the modular finite quotient
is

$$
C=\SL_2(\F_3),\quad
 \Gamma=\widehat\Z^\times,\quad
 \theta_\gamma=\Int(\operatorname{diag}(-1,1))
       \text{ when }\gamma\bmod3=-1.
$$

The finite quotient is $\GL_2(\F_3)$ of order48,
and the complementary order-two subgroup is not normal.
A separate central-intersection issue already occurs for
a cyclic order-four cusp in $\operatorname{Dic}_3$.
\end{statement}
\begin{proof}
Let $J=\operatorname{diag}(-1,1)$ and
$c=\left(\begin{smallmatrix}1&1\\0&1\end{smallmatrix}\right)$
over $\F_3$. Then $JcJ^{-1}=c^{-1}\ne c$.
In $C\rtimes C_2$, conjugating the complement generator by
$c$ gives a nontrivial $C$-coordinate, so leaves the complement.
The determinant sequence identifies this semidirect product
with $\GL_2(\F_3)$.
The modular subgroup generated by the two squared cusp
generators has full image $\SL_2(\F_3)$: their reductions
are upper and lower elementary matrices with nonzero
parameters. Thus triviality of the induced action is a real
condition, not a consequence of the modular relations.

More generally, when $C_0\triangleleft C$ is stable and the
relevant cusp is already contained in $C_0$, the subgroup
$C_0\rtimes\Gamma$ is normal exactly when $\Gamma$ acts
trivially on $C/C_0$. This follows by computing
$(c,1)(1,\gamma)(c^{-1},1)$.
If the cusp is not contained in $C_0$, its actual intersection
must instead be retained in the diagonal quotient.
For $\operatorname{Dic}_3=C_3\rtimes C_4$, an order-four
cusp has central square. Its inner image has order two,
while its intersection with the trivial subgroup is trivial.
Replacing one of those kernels by the other therefore
changes the reconstruction.
\end{proof}

\begin{statement}{B.19}{Proved}{The master fibre-product exactness calculation}
Let $A\to\Gamma$, $B\to\Gamma$, and $G\to\Gamma$ be
homomorphisms. Their iterated fibre product is

$$
P=\{(a,b,g):\alpha(a)=\beta(b)=\gamma(g)\}.
$$

If $A\to\Gamma$ and $B\to\Gamma$ are epimorphisms in
the category in use, then $P\to G$ is an epimorphism with
kernel $\ker\alpha\times\ker\beta$.
Without those hypotheses its image is exactly

$$
\gamma^{-1}(\im\alpha\cap\im\beta).
$$

\end{statement}
\begin{proof}
The kernel formula follows by setting $g=1$.
A point has a lift precisely when its common quotient value
lifts to both $A$ and $B$. For sheaves the lifts may be
chosen on a common refinement of the two covers, giving
the same sheaf-image formula.
This proves the formal exact sequences used in the
fibre-product constructions around C3-559--562 and
C3-673--681 whenever the stated epimorphism hypotheses hold.
It also identifies exactly what must be checked before
asserting their surjectivity for restricted subgroups.
\end{proof}

\begin{statement}{B.20}{Proved}{The affine-cocycle obstruction to a diagonal section}
Let $N$ map to an affine group $A\rtimes\Gamma$, with
translation coordinate $\ell$ and a normal subgroup $L$
mapping to translations $jL$. The induced translation
class on the quotient takes values in $A/jL$.
The image of the diagonal subgroup is exactly the zero
locus of this class. A prescribed diagonal section exists
only if this class is zero.
Even with trivial linear action, the extension can fail
to have any section.
\end{statement}
\begin{proof}
An affine element can be made diagonal by multiplication
with an allowed translation precisely when its translation
coordinate belongs to $jL$. This is the stated condition
in the quotient module.
For an explicit nonsplit example use

$$
0\longrightarrow2\widehat\Z\longrightarrow\widehat\Z
 \longrightarrow C_2\longrightarrow0,
$$

represented by translations of the affine line.
The nontrivial quotient class has odd translation coordinate.
A section would require an odd element of order two in
the torsion-free group $\widehat\Z$, which is impossible.
The same obstruction is already visible in
$0\to2(\Z/4)\to\Z/4\to C_2\to0$.
Thus a Borel realization does not automatically supply the
particular diagonal splitting requested by a later diagram.
\end{proof}

\begin{statement}{B.21}{Proved}{Restricted functors and their exact images}
An automorphism of $C$ descends to $C/K$ exactly when it
preserves $K$ as a subgroup, not merely as an unnamed
isomorphism type. If $A'\subseteq A$ and $B'\subseteq B$
are restrictions in B.19, their fibre-product projection
has image

$$
\gamma^{-1}(\im A'\cap\im B')
$$

and kernel $\ker(A'\to\Gamma)\times\ker(B'\to\Gamma)$.
These formulas do not establish that the source-specific
restricted projections on C3-690 are epimorphisms.
\end{statement}
\begin{proof}
A quotient map factors through $C/K$ exactly when it kills
$K$. Applying the same assertion to the inverse automorphism
gives equality of the preserved subgroup.
The fibre-product image and kernel follow from B.19 with
the restricted domains. In the profinite case one retains
closed subgroups and continuous maps; in the sheaf case one
uses sheaf images. The missing source-specific assertion
is precisely equality of these images with the required
quotient image. No such equality is inferred from a notation
for the fibre product.
\end{proof}

\begin{statement}{B.22}{Proved}{Diagonal gluing of marked extensions}
Let $G_i\to\Gamma$ be extensions with the same marked kernel
$C$, and let $P=\prod_\Gamma G_i$.
A subgroup $H\subseteq P$ gives simultaneous compatible
identifications of the extensions exactly when

$$
H\to\Gamma\text{ is an epimorphism},\qquad
 H\cap C^I=\Delta C.
$$

Then every projection $H\to G_i$ is an isomorphism.
\end{statement}
\begin{proof}
The kernel of a projection lies in $\Delta C$ and has
identity $i$th coordinate, so is trivial.
Given $g_i$, lift its quotient value to an element of $H$.
Their difference in the $i$th coordinate belongs to $C$;
multiplication by that diagonal element corrects the
coordinate. This proves surjectivity, locally in the
sheaf setting. Conversely compatible identifications have
a graph with precisely these properties.
The transition identifications obey the cocycle relation
because they arise from projections of one common group.
\end{proof}

\begin{statement}{B.23}{Proved}{The finer parameter scheme on C4-41}
Use the earlier visible convention $\nu=2h+i$ from C4-38,
with constant bits $i,j\in\{0,1\}$.
Put

$$
R=\begin{pmatrix}0&-1\\1&-1\end{pmatrix},\qquad
 S=\begin{pmatrix}0&1\\1&0\end{pmatrix},\quad
 \alpha=(a,b),\quad\beta=(c,d).
$$

The equations on C4-39--41 become

$$
a+b-c+d=0,\qquad 3b-2h=i+j,\qquad
 4h+2i+1\text{ invertible}.
$$

Each bit chart is

$$
\Spec\Z[c,d,w,(1-12w+6i+4j)^{-1}],
 \qquad w=b-h.
$$

The complete scheme is therefore smooth and flat of relative
dimension three over all of $\Z$, including two and three.
Its fibres at those two primes are four affine three-spaces.
The bits here are constant components, not $\mu_2$ parameters.
\end{statement}
\begin{proof}
Put $k=i+j$. The integral change of variables is

$$
b=k-2w,\qquad h=k-3w,\qquad
 a=c-d-k+2w.
$$

It is inverse to $w=b-h$ on the defining equations.
The invertible multiplier becomes
$1-12w+6i+4j$.
This is an integral unimodular calculation; a saturated
Groebner basis over $\Q$ would not suffice.
At two the multiplier is one, and at three it is $1+j$,
always a unit. Thus all four bad-prime charts are affine
three-spaces. For primes at least five, one value of $w$
is excluded in each chart, giving $4p^2(p-1)$ points.

The point counts, independently enumerated, are
$32,108,400,1176,4840,8112$ at
$p=2,3,5,7,11,13$, and $256,2048,2916$ over
$\Z/4,\Z/8,\Z/9$.
Over $\Z$, the only possible bit and $w$ values are
$(i,j,w)=(0,0,0)$ and $(1,1,1)$, with $c,d$ arbitrary,
because an integral unit multiplier is $1$ or $-1$.

For completeness replace the right side by
$2h+s_2i+s_3j$, $s_2,s_3\in\{\pm1\}$.
The same calculation gives multiplier
$1-12w+6s_2i+4s_3j$.
All charts remain smooth. If $s_3=-1,j=1$, their fibre
at three is empty; this is a substantive change.
Changing the convention for the $i$-sign is instead
absorbed by $h_-=h_++i$, $w_-=w_+-i$.
The unaugmented linear map in $(a,b,c,d)$ has Smith
invariants $1,3$; the coefficient two of $h$ removes that
cokernel. This explains why failure of that smaller map
to be onto at three does not make the finer scheme singular.
\end{proof}
